Một trí thức chính trực vừa nằm xuống. Mọi người nhớ ra rằng ông đã phản biện bao điều ngay thẳng, từ nhiều năm trước, tiếc là không ai có trách nhiệm tiếp thu.

Chẳng hạn, ông nói với Thủ tướng từ đầu thập niên 1960 rằng, hệ thống kinh tế Xô Viết sẽ sụp đổ nếu tiếp tục vận hành thế này. Người ta trầm trồ khen ông nhìn xa. Hiện nay, năng suất lao động của Việt Nam còn thấp hơn Campuchia; vậy thì câu nói của ông ngày nay phải hiểu thế nào?

Ví như, năm 1992 trong một góp ý cho Hiến pháp ở Uỷ ban Trung ương MTTQ Việt Nam, ông đề nghị “tạm gác lại” việc đưa vào Hiến pháp các thuật ngữ “chủ nghĩa xã hội, chủ nghĩa Marx-Lenin”, “mô hình CNXH, cũng như Marx-Lenin không còn đáp ứng được mục tiêu phát triển đất nước trong giai đoạn hiện nay”.

Hệ quả là, ông cho rằng phải có “nhà nước pháp quyền”, chứ không thể tiếp tục “chuyên chính vô sản” hay “pháp quyền nửa vời”.

Mấy hôm nay, báo chí và dư luận ca ngợi ông. Người ta quên rằng có lúc ông từng bị quy chụp. Vài hôm nữa, khi ông hoàn toàn yên nghỉ, mọi việc lại như xưa. 

Lại có vài trí thức cương trực phản biện thẳng thắn, mà không ai nghe. Đôi lúc họ bị chụp mũ. Đến khi có người trong số chính trực ấy nằm xuống, thiên hạ trong nỗi tiếc thương lại ồn ào nhớ ra rằng ông ấy hay bà ấy đã từng góp ý rất sáng suốt, từ nhiều năm trước.

Giá như thiên hạ chịu lắng nghe một vài điều mà những người chính trực phản biện, thì xã hội bớt tang thương. Người chính trực lúc ra đi bớt ngậm ngùi. Người ở lại tránh được tiếng đạo đức giả. Còn xã hội tránh được những cuộc nhỡ tầu triền miên.   

Người chính trực đương nhiên không cần được ngợi ca, một khi phản biện của họ bị vứt xó. 
Hà Nội, 15/5/2018